% ============================================================
%  Trolley Retrieval -- main writeup
%  Compile with: lualatex writeup.tex  (twice for refs)
% ============================================================
\documentclass{article}

\usepackage[final,nonatbib,main]{neurips_2026}  % final+main = non-anonymous, no line nums

% suppress NeurIPS conference footer notice
\makeatletter
\renewcommand{\@notice}{}
\makeatother

\usepackage[T1]{fontenc}
\usepackage{hyperref}
\usepackage{url}
\usepackage{booktabs}
\usepackage{amsfonts}
\usepackage{amsmath}
\usepackage{amssymb}
\usepackage{amsthm}
\usepackage{nicefrac}
\usepackage{microtype}
\usepackage{xcolor}
\usepackage{array}
\usepackage{tabularx}
\usepackage{mathtools}
\usepackage{enumitem}

\hypersetup{
  colorlinks=true,
  linkcolor=blue!60!black,
  citecolor=green!50!black,
  urlcolor=blue!60!black
}

% -- theorem environments ------------------------------------------------------
\newtheorem{theorem}{Theorem}[section]
\newtheorem{lemma}[theorem]{Lemma}
\newtheorem{corollary}[theorem]{Corollary}
\newtheorem{proposition}[theorem]{Proposition}
\newtheorem{definition}[theorem]{Definition}
\newtheorem{remark}[theorem]{Remark}
\newtheorem{example}[theorem]{Example}
\newtheorem{claim}[theorem]{Claim}

% -- macros --------------------------------------------------------------------
\newcommand{\N}{\mathbb{N}}
\newcommand{\Np}{\mathbb{N}^{+}}
\newcommand{\Z}{\mathbb{Z}}
\newcommand{\Yes}{\textsf{Yes}}
\newcommand{\No}{\textsf{No}}
\DeclareMathOperator{\sgn}{sgn}

% small breakable typewriter for long Lean identifiers; "." and "_" become
% break opportunities (use \us for underscore, \dt for dot inside \lean).
\newcommand{\us}{\discretionary{}{}{}\_\discretionary{}{}{}}
\newcommand{\dt}{.\discretionary{}{}{}}
\newcommand{\lean}[1]{{\small\ttfamily #1}}

% -- document ------------------------------------------------------------------
\title{Trolley Retrieval:\\
A Complete Analysis with a Machine-Checked\\
Distinguishability Theorem}

\author{%
  Akaash Dash \\
  \texttt{akaash.dash@gmail.com}
}

\begin{document}

\maketitle

% ============================================================
\begin{abstract}
We analyse the \emph{Trolley Retrieval} problem (G-Research GRiddles
Series~A, Puzzle~4). A trolley sits on an unknown cell of a one-way
infinite tape whose cell~$n$ carries the label~$c(n)$, where
$c$ is a permutation of the positive integers known to the player Bob.
Each turn Bob pushes the trolley one cell left or right and is told
whether the label of the new cell is strictly less than the current
time. Bob must eventually name the trolley's exact position; falling off
the left end, or guessing wrongly, loses.

We prove the answer: \textbf{Bob has a winning strategy if and only if
$c$ is \emph{not} eventually the identity}, i.e.\ $c(n)\neq n$ for
infinitely many~$n$ (equivalently, $c$ is not a finitary permutation).
We give the full elementary proof of both directions. The losing
direction rests on a parity lock: the threshold Bob can probe is always
\emph{even}, so an identity tail makes the pair $(2j,\,2j{+}1)$
permanently indistinguishable. The winning direction rests on a
\emph{distinguishability theorem}: for a non-eventually-identity~$c$, no
pair of starting cells is indistinguishable. Its crux---previously
believed to require global surjectivity bookkeeping---falls to an
elementary forward-closure argument.

The losing direction, the distinguishability theorem, and explicit
winning strategies for two natural ``locally decodable'' families are
formalised and machine-checked in Lean~4 and are axiom-clean. We are
scrupulous about the boundary of what is proven: a single, uniform,
explicitly formalisable winning strategy for \emph{every}
non-eventually-identity~$c$ remains open, owing to a genuine
timing/resolution tension that we describe precisely.
\end{abstract}

% ============================================================
\section{The problem}\label{sec:problem}

The exact statement is as follows. Let $c(1), c(2), \dots$ be a sequence
of positive integers in which every positive integer appears exactly
once; equivalently $c$ is a permutation (bijection) of $\Np$. A one-way
infinite sequence of cells is given: cell~$n$ (the $n$-th from the left,
$n\ge 1$) is labelled~$c(n)$.

At time $t=0$ an invisible trolley is placed on an unknown cell, the
$k_0$-th from the left. For each integer $t\ge 1$, in order:
\begin{enumerate}[topsep=2pt,itemsep=1pt]
  \item Bob chooses a direction, left or right; the trolley moves one
  cell that way, to cell $k_t$.
  \item The trolley sends \Yes{} if $c(k_t) < t$, and \No{} otherwise.
\end{enumerate}
If ever $k_t = 0$ (the trolley falls off the left end), Bob loses
immediately. At any time Bob may instead \emph{guess} the trolley's
position; a correct guess retrieves it, an incorrect guess loses.

\begin{quote}
\itshape Determine all permutations $c$ for which Bob has a strategy
guaranteeing safe retrieval.
\end{quote}

Bob knows $c$ in full; the only hidden datum is the start cell $k_0$. The
trolley is invisible---Bob never sees $k_t$ directly. His sole feedback
is the binary signal, and he chooses each move adaptively from the
signal history.

% ------------------------------------------------------------
\subsection{The answer}

\begin{theorem}[Main]\label{thm:main}
Bob has a strategy ensuring safe retrieval if and only if $c$ is
\emph{not eventually the identity}; that is, if and only if
$c(n)\neq n$ for infinitely many~$n$.
\end{theorem}

We write $c$ is \emph{eventually identity} (EI) if there is a threshold
$N$ with $c(n)=n$ for all $n>N$. Since $c$ is a bijection, ``not EI'' has
several equivalent forms, all used below:
\begin{itemize}[topsep=2pt,itemsep=1pt]
  \item $\{\,n : c(n)\neq n\,\}$ is infinite ($c$ is not finitary);
  \item $c^{-1}$ is not eventually identity (the condition is symmetric
  under $c\leftrightarrow c^{-1}$);
  \item $c$ has infinitely many \emph{upward gaps}: positions $p$ with
  $c(p{+}1)-c(p)\ge 2$.
\end{itemize}
The last equivalence is worth recording, as it is exactly what fails in
the EI case.

\begin{proposition}\label{prop:gaps}
$c$ has infinitely many upward gaps if and only if $c$ is not eventually
identity.
\end{proposition}
\begin{proof}
If $c(n{+}1)\le c(n)+1$ for all $n>N$ (no gaps past $N$), then on
$(N,\infty)$ the bijection $c$ is non-decreasing by steps of at most~$1$;
a strictly value-preserving bijection of a tail of $\Np$ that never
jumps up by~$2$ cannot skip a value and cannot repeat one, forcing
$c(n)=n$ for all $n>N$. The contrapositive is the claim.
\end{proof}

% ------------------------------------------------------------
\subsection{Reformulation: the alarm-clock view and the parity lock}
\label{sec:reform}

It is convenient to read the labels as wake-up times. Think of cell~$n$
as carrying an alarm clock set for time $c(n)$: at every time
$t > c(n)$ the cell is ``awake'', and a trolley standing on it reports
\Yes; while $t \le c(n)$ the cell is ``asleep'' and reports \No. Because
$c$ is a bijection, exactly one cell wakes at each integer time, and
every cell wakes eventually. This is just a restatement of
$\text{signal}=\Yes \iff c(k_t) < t$.

Let $D_t := k_t - k_0$ be Bob's net displacement at time~$t$.

\begin{lemma}[Parity lock]\label{lem:parity}
As long as Bob has not yet guessed, $D_t \equiv t \pmod 2$. Equivalently,
$t - D_t$ is always even.
\end{lemma}
\begin{proof}
$D_0=0$, and each move changes $D$ by $\pm 1$ and $t$ by $+1$, so the
parity of $t-D_t$ is invariant; it starts even.
\end{proof}

The parity lock is the single structural fact behind the losing
direction. A second, equally elementary identity exposes \emph{why}.
Let $L_t$ be the number of \emph{left} moves Bob has made by time~$t$ and
$R_t$ the number of right moves, so $t=L_t+R_t$ and $D_t=R_t-L_t$. Then
\begin{equation}\label{eq:tD}
  t - D_t \;=\; (L_t+R_t) - (R_t-L_t) \;=\; 2L_t .
\end{equation}
Write $g(m) := c(m)-m$ for the displacement of the label from the
diagonal. The signal at time~$t$ is
\begin{align}
  \text{\Yes}
  &\iff c(k_0 + D_t) < t
  \iff g(k_0+D_t) < t - D_t - k_0 \notag\\
  &\iff g(k_0+D_t) < 2L_t - k_0 . \label{eq:Lthreshold}
\end{align}
In other words, the walk reads the sequence~$g$ against a threshold
$2L_t - k_0$ that Bob controls only through his \emph{left-move count}.

\begin{remark}[The whole obstruction, in one line]\label{rem:flat}
For the identity, $g\equiv 0$, so~\eqref{eq:Lthreshold} collapses to
$\Yes \iff k_0 < 2L_t$. Bob can only ever compare $k_0$ against
\emph{even} thresholds, and $2j < 2L \iff 2j{+}1 < 2L$ for every~$L$.
Hence the start cells $2j$ and $2j{+}1$ produce \emph{identical signals
under every strategy}: the parity bit of $k_0$ is unrecoverable. This is
the seed of the losing direction; the work below is to extend it to every
eventually-identity~$c$ despite the perturbed initial segment.
\end{remark}

% ============================================================
\section{Losing direction: eventually identity $\Rightarrow$ Bob loses}
\label{sec:losing}

\begin{theorem}\label{thm:losing}
If $c$ is eventually identity, Bob has no strategy that wins from every
start cell.
\end{theorem}

The proof isolates a structural ``parallel-evolution'' lemma and then
verifies its hypothesis for two carefully chosen start cells.

% ------------------------------------------------------------
\subsection{Parallel evolution}

Fix a strategy~$\sigma$. For a start cell $k$ let $\mathrm{pos}(k,t)$ be
the position and $\mathrm{hist}(k,t)$ the signal history (most-recent
first) at time~$t$. Bob's action at any decision point is a function of
$\mathrm{hist}$ alone.

\begin{lemma}[Parallel evolution]\label{lem:parallel}
Let $k_1, k_2$ be two start cells. Suppose that, started from~$k_1$, Bob
makes no guess before time~$t$, and that at every time $s\le t$ the
signal observed from $k_1$'s trajectory equals the signal observed from
$k_2$'s trajectory. Then the two histories agree at time~$t$, and the
positions stay rigidly offset:
\[
  \mathrm{hist}(k_1,t) = \mathrm{hist}(k_2,t),
  \qquad
  \mathrm{pos}(k_2,t) = \mathrm{pos}(k_1,t) + (k_2-k_1).
\]
\end{lemma}
\begin{proof}
Induction on~$t$. At $t=0$ the histories are empty and the offset is
$k_2-k_1$ by definition. For the step, equal histories make $\sigma$
choose the \emph{same} action in both worlds; since Bob has not guessed,
that action is a move, applied identically to both, preserving the
offset. The new signal entries are equal by the time-$(t{+}1)$ hypothesis,
so the histories remain equal.
\end{proof}

\begin{corollary}\label{cor:notboth}
If $k_1\neq k_2$ and the two signal traces agree at \emph{every} time
under~$\sigma$, then $\sigma$ cannot win from both $k_1$ and $k_2$.
\end{corollary}
\begin{proof}
Equal traces force equal histories at all times (Lemma~\ref{lem:parallel}
applied up to either guess time), hence equal first-guess times
$T_1=T_2=:T$ and equal guesses $g_1=g_2$. But correctness would require
$g_1 = \mathrm{pos}(k_1,T)$ and $g_2=\mathrm{pos}(k_2,T)$, while
$\mathrm{pos}(k_2,T)=\mathrm{pos}(k_1,T)+(k_2-k_1)$. These contradict
$k_1\neq k_2$.
\end{proof}

% ------------------------------------------------------------
\subsection{The indistinguishable pair for eventually-identity $c$}

Suppose $c(n)=n$ for all $n>N$. Let $M$ be an upper bound for the finitely
many perturbed labels: $M \ge \max\{\,c(m) : 1\le m\le N\,\}$ and $M\ge
N+1$. Choose $j$ with $2j > M$ and set
\[
  k_1 = 2j, \qquad k_2 = 2j+1 .
\]

\begin{lemma}[Signal match]\label{lem:sigmatch}
Fix a no-guess strategy run from $k_1=2j$. At every time~$t$, every
position~$p\ge 1$ reachable from $2j$ in $t$ moves (so $2j - t \le p$)
that respects the parity lock ($t+p$ even) satisfies
\[
  \text{signal at } p \text{ at time } t \;=\;
  \text{signal at } p{+}1 \text{ at time } t .
\]
\end{lemma}
\begin{proof}
Two cases on the location~$p$.

\emph{(a) $p > N$ (identity region).} Both $p$ and $p{+}1$ exceed~$N$, so
$c(p)=p$ and $c(p{+}1)=p{+}1$. The signal at $p$ is $\Yes \iff p<t$, and
at $p{+}1$ is $\Yes\iff p{+}1<t$. These differ only if $t=p+1$. But then
$t+p = 2p+1$ is odd, violating the parity hypothesis $t+p$ even. So the
signals agree.

\emph{(b) $p \le N$ (perturbed region).} Here the trolley has travelled
from $2j$ to $p\le N$, so $t \ge 2j - p \ge 2j - N > M - N \ge \dots$;
concretely $2j>M\ge N+1$ forces $t > M$. Both $c(p)$ and $c(p{+}1)$ are at
most~$M$ (using $c(N{+}1)=N{+}1\le M$ for the boundary), hence
$c(p), c(p{+}1) < t$: both signals are \Yes. They agree.
\end{proof}

The two cases capture the dichotomy of Remark~\ref{rem:flat}: in the
identity tail, the parity lock makes $2j$ and $2j{+}1$ collide; in the
perturbed segment, the time cost of reaching it ($t\ge 2j-N$) is so large
that every cell there is already awake, so both worlds report \Yes{} and
again collide.

\begin{proof}[Proof of Theorem~\ref{thm:losing}]
Suppose $\sigma$ wins from every start cell, in particular from
$k_1=2j$ and $k_2=2j{+}1$ above. We claim the two signal traces agree at
every time, which contradicts Corollary~\ref{cor:notboth}.

By induction on~$t$, the trajectories from $2j$ and $2j{+}1$ evolve in
lockstep with positions offset by exactly~$1$ and equal histories, as
long as Bob has not guessed: the inductive step applies the signal match
(Lemma~\ref{lem:sigmatch}) at $\mathrm{pos}(2j,t{+}1)$, whose
hypotheses---validity $p\ge 1$ (Bob wins, so the trajectory is safe),
reachability $2j-t\le p$ (a $\pm1$ walk), and parity (Lemma~\ref{lem:parity})---all
hold. Hence the signal entering both histories at each step is identical,
so $\mathrm{hist}(2j,t)=\mathrm{hist}(2j{+}1,t)$ and
$\mathrm{pos}(2j{+}1,t)=\mathrm{pos}(2j,t)+1$ for all~$t$. The traces
agree at every time, and Corollary~\ref{cor:notboth} gives the
contradiction.
\end{proof}

\begin{remark}[The half-infinite tape is irrelevant here]
For the chosen large $k_0\approx 2j$, no finite-time strategy can reach
cell~$0$, so the ``falling off'' rule never binds; the indistinguishability
argument is unaffected by the absence of a left end.
\end{remark}

% ============================================================
\section{Winning direction: not eventually identity $\Rightarrow$ Bob wins}
\label{sec:winning}

The heart of the winning direction is a converse to the losing
mechanism: when $c$ is not eventually identity, \emph{no} pair of start
cells can stay indistinguishable. We make ``indistinguishable'' precise,
prove the theorem, and then explain how distinguishability is turned into
a winning strategy---fully for two natural families, and with an honest
account of the general gap.

% ------------------------------------------------------------
\subsection{The distinguishability theorem}

Reaching displacement~$D$ from the start costs time at least~$|D|$, and
the parity lock forces $t\equiv D\pmod 2$. Within those constraints Bob
can pad with oscillation to reach any admissible time, and (because a cell
once awake stays awake) the only information at displacement~$D$ is the
\emph{set of times $t\ge D$, $t\equiv D\pmod 2$, at which the cell is
awake}---equivalently, the single comparison $c(a+D) < t$. This motivates:

\begin{definition}[Indistinguishable pair]\label{def:indist}
Two start cells $a<b$ are \emph{indistinguishable} (for~$c$) if for every
displacement $D\ge 0$ and every time $t\ge D$ with $t\equiv D\pmod 2$,
\[
  c(a+D) < t \iff c(b+D) < t .
\]
\end{definition}

Quantifying over \emph{all} reachable $(D,t)$---including small times that
catch early wake-ups---is essential; it is exactly the freedom that lets
Bob's oscillation observe a cell's wake-up time at full resolution.

\begin{theorem}[Distinguishability]\label{thm:lemmaA}
If $a<b$ are indistinguishable for~$c$, then $b=a+1$ and $c(n)=n$ for all
$n\ge a$. Consequently, if $c$ is not eventually identity, then \emph{no}
pair of start cells is indistinguishable.
\end{theorem}

We prove the first sentence; the second is its contrapositive (an
indistinguishable pair would exhibit an identity tail, making $c$ EI).
Throughout write $\Delta := b-a \ge 1$.

The proof has two regimes by the size of~$\Delta$. The case $\Delta=1$ is
a clean shift argument; the case $\Delta\ge 2$ is the crux. We first
extract the local structure shared by both.

\begin{lemma}[Consecutiveness]\label{lem:consec}
Call a displacement~$D$ \emph{active} (for the cell~$a$) if
$c(a+D)\ge D$, and similarly for~$b$. If $a<b$ are indistinguishable and
$D$ is active for both $a$ and $b$, then $c(a+D)$ and $c(b+D)$ are
consecutive integers:
\[
  |\,c(a+D) - c(b+D)\,| = 1 .
\]
\end{lemma}
\begin{proof}
Let $X=c(a+D)$, $Y=c(b+D)$. Injectivity gives $X\neq Y$; say $X<Y$. If
$Y\ge X+2$, pick~$t$ with $X<t\le Y$ of the parity $t\equiv D\pmod 2$
(possible because the interval $(X,Y]$ has length $\ge 2$ and so contains
both parities). Then $c(a+D)<t$ but $c(b+D)\not<t$, with $t\ge D$ (as $D$
is active and $t>X\ge D$), contradicting indistinguishability at $(D,t)$.
Hence $Y=X+1$. The case $Y<X$ is symmetric.
\end{proof}

\begin{lemma}[Bad displacements are forward-closed]\label{lem:badprop}
Call $D$ \emph{bad} (for~$a$) if $c(a+D) < D$. If $a<b$ are
indistinguishable and $D$ is bad, then $D+\Delta$ is bad: $c(a+D+\Delta) < D+\Delta$.
\end{lemma}
\begin{proof}
$D$ bad means $c(a+D)<D$, i.e.\ the cell $a+D$ is awake at the admissible
time $t$ with $t\equiv D \pmod 2$, $D\le t$, as soon as $t>c(a+D)$; take
the least such admissible~$t$, so $t\le D+1$ and the signal is \Yes.
Indistinguishability transfers \Yes{} to $b+D = a+(D+\Delta)$:
$c(a+D+\Delta) < t \le D+1$, and since the value is an integer
$c(a+D+\Delta)\le D < D+\Delta$. (The point is purely that an
already-saturated \Yes{} propagates by $+\Delta$; this is where the
constraint $b=a+\Delta$ enters.)
\end{proof}

Lemma~\ref{lem:badprop} is decisive: on each residue class modulo
$\Delta$, the bad displacements form an \emph{up-set}, so the
\emph{active} (non-bad) displacements form a finite or cofinite
\emph{prefix}---never an interleaving of good and bad. This is exactly the
structural fact that earlier, more elaborate attempts (parity arguments,
value-hogging bookkeeping) were trying to recover. With it, the crux is
short.

\begin{lemma}[Chain bound]\label{lem:chain}
If $a<b$ are indistinguishable, fix a residue $r$ with $0\le r<\Delta$.
For any~$k$ such that $r, r+\Delta, \dots, r+k\Delta$ are all active,
\[
  c(a + r + k\Delta) \;\le\; c(a+r) + k .
\]
\end{lemma}
\begin{proof}
By Lemma~\ref{lem:consec}, consecutive active displacements on the
residue satisfy $|c(a + r + (i{+}1)\Delta) - c(a + r + i\Delta)| = 1$
(apply consecutiveness at the displacement $r+i\Delta$, whose
$+\Delta$ shift lands at $b + r + i\Delta$). Summing $k$ unit steps and
the triangle inequality give $c(a+r+k\Delta)\le c(a+r)+k$.
\end{proof}

\begin{lemma}[$\Delta\ge 2$ is impossible]\label{lem:Delta2}
If $a<b$ are indistinguishable with $\Delta=b-a\ge 2$, contradiction.
\end{lemma}
\begin{proof}
Let $M := \max_{0\le r<\Delta} c(a+r)$. Suppose a displacement $D$ is
active, and write $D=r+k\Delta$ with $0\le r<\Delta$. By forward-closure
(Lemma~\ref{lem:badprop}), \emph{all} of $r, r+\Delta, \dots, r+k\Delta$
are active (active is a prefix on each residue). Activity gives
$c(a+D)\ge D = r+k\Delta$, while the chain bound gives
$c(a+D)\le c(a+r)+k \le M+k$. Hence
\[
  r + k\Delta \le M + k
  \quad\Longrightarrow\quad
  k(\Delta-1) \le M - r \le M .
\]
With $\Delta\ge 2$ this bounds $k\le M$, so $D = r+k\Delta < (M{+}1)\Delta$.
Thus \emph{every} displacement $D \ge (M{+}1)\Delta$ is bad:
$c(a+D) < D$ for all such~$D$.

This forces a pigeonhole collapse. For $D \ge (M{+}1)\Delta$, the cell at
position $a+D$ has label $c(a+D) < D = (a+D) - a$, i.e.\ $c(p) < p-a$ for
all large~$p$. The finitely many small positions have labels bounded by
some constant. Choosing $N$ large enough, $c$ maps the segment $[1,N]$
into $[1, N-a-1]$: large positions $p$ land below $p-a \le N-a-1$, and
the bounded small positions also fit once $N$ is large. But $[1,N]$ has
$N$ elements and $[1, N-a-1]$ has $N-a-1 < N$, so $c$ cannot be injective
on $[1,N]$. Contradiction.
\end{proof}

\begin{remark}[Why this is the right proof]
The contradiction in Lemma~\ref{lem:Delta2} uses \emph{no} parity
argument and \emph{no} even/odd split on~$\Delta$. The single lever is
forward-closure of badness (Lemma~\ref{lem:badprop}), which turns the
``active'' set on each residue into a prefix; the chain bound then makes
that prefix finite for $\Delta\ge 2$, and a one-line pigeonhole finishes.
This is markedly simpler than---and subsumes---the earlier even/odd
case analysis.
\end{remark}

It remains to handle $\Delta=1$.

\begin{lemma}[$\Delta=1$ forces an identity tail]\label{lem:Delta1}
If $a<b=a+1$ are indistinguishable, then $c(n)=n$ for all $n\ge a$.
\end{lemma}
\begin{proof}
For $\Delta=1$, $b+D = a+(D+1)$, so consecutiveness
(Lemma~\ref{lem:consec}) on active displacements reads
$|c(a+D+1)-c(a+D)|=1$: the labels along $a, a{+}1, a{+}2, \dots$ form a
self-avoiding $\pm 1$ walk. Let $\sigma(0)=\sgn\bigl(c(a{+}1)-c(a)\bigr)$
be the sign of the first step. If $\sigma(0)=-1$ (descending start), the
walk $c(a), c(a){-}1, c(a){-}2,\dots$ would eventually drop below~$1$,
impossible for labels in $\Np$; so $\sigma(0)=+1$. A self-avoiding $\pm 1$
walk that starts ascending must ascend forever (any later descent would
revisit a value, breaking injectivity), giving $c(n+1)=c(n)+1$ for all
$n\ge a$, i.e.\ $c(n)=n + C$ on $[a,\infty)$ with $C=c(a)-a$. If $C>0$,
the values $\{1,\dots,C\}$ are absent from $c([a,\infty))$ and must be
supplied by the finite set $c([1,a{-}1])$, which has only $a-1$ elements;
bijectivity of $c$ on $\Np$ then forces $C=0$ (and $C<0$ was already
excluded). Hence $c(n)=n$ for $n\ge a$. (In Lean this sign analysis is the
pair \lean{indist\us Delta1\us pos} / \lean{indist\us Delta1\us neg}.)
\end{proof}

\begin{proof}[Proof of Theorem~\ref{thm:lemmaA}]
Let $a<b$ be indistinguishable, $\Delta=b-a$. By Lemma~\ref{lem:Delta2},
$\Delta\ge 2$ is impossible, so $\Delta=1$, i.e.\ $b=a+1$; then
Lemma~\ref{lem:Delta1} gives $c(n)=n$ for $n\ge a$, so $c$ is eventually
identity. Contrapositively, if $c$ is not eventually identity then no
pair is indistinguishable.
\end{proof}

% ------------------------------------------------------------
\subsection{From distinguishability to a strategy}
\label{sec:strategy}

Theorem~\ref{thm:lemmaA} says that for a non-EI~$c$ every pair of cells
admits a separating observation. A winning strategy must (i) narrow the
start cell down to a small candidate set, then (ii) reach a separating
observation while keeping the trolley on the tape. Safety is free if Bob
only ever oscillates and moves right: positions stay $\ge k_0 \ge 1$, so
the trolley never falls off.

\paragraph{Phase~1 (bracket to two candidates).}
Bob oscillates the displacement between $0$ and~$1$, moving
\emph{right first}, so $D_t\in\{0,1\}$ and $k_t\ge k_0$. At every even
time $t=2m$ the trolley sits on $k_0$, with signal $\Yes \iff c(k_0)<2m$.
Since $c(k_0)$ is finite, there is a first even time $T_1=2m^\ast$ with a
\Yes, pinning $c(k_0)\in\{T_1{-}2,\,T_1{-}1\}$, hence
\[
  k_0 \in \{\,c^{-1}(T_1{-}2),\; c^{-1}(T_1{-}1)\,\}
\]
(a singleton if one of those preimages is~$0$, in which case Bob guesses
at once). So Phase~1 always narrows $k_0$ to (at most) two candidates,
the cells waking at one even and the next odd time.

\paragraph{Phase~2 (resolve).}
Bob must now separate the two candidates. By Theorem~\ref{thm:lemmaA} the
candidate pair is distinguishable: some admissible $(D,t)$ with $t\ge D$,
$t\equiv D\pmod2$ yields different signals in the two worlds. Reaching it
and reading the signal identifies $k_0$ and hence the current position,
which Bob guesses.

The subtlety, and the location of the only remaining gap, is in Phase~2:
\emph{which} separating observation to aim for, and how to monitor it at
the necessary time resolution. We make this precise in
Section~\ref{sec:status} after first giving complete strategies for two
natural families where Phase~2 is finite and local.

% ------------------------------------------------------------
\subsection{Two complete winning families (locally decodable)}
\label{sec:local}

When the candidate pair is always separated by a wake-up that occurs at a
\emph{bounded} displacement, Bob can resolve it by a fixed-amplitude
oscillation that monitors a small window from time~$0$---catching even the
early wake-ups that a navigate-then-observe approach would arrive too late
to see. Two regimes are formalised.

\begin{definition}[Locally $K{=}2$ decodable]\label{def:k2}
$c$ is \emph{locally decodable} if the map
$k_0 \mapsto \bigl(\text{first even \Yes{} time of } c(k_0),\;
\text{first odd \Yes{} time of } c(k_0{+}1)\bigr)$
is injective; equivalently the bucketed pair $(c(k_0), c(k_0{+}1))$
determines~$k_0$.
\end{definition}

\begin{theorem}[$K{=}2$ winning theorem]\label{thm:k2}
If $c$ is locally decodable, Bob wins. The strategy oscillates
$D\in\{0,1\}$, reads off $c(k_0)$ and $c(k_0{+}1)$ from the first even
and first odd \Yes, decodes~$k_0$, and guesses its current
position---never leaving $\{k_0, k_0{+}1\}$, so safety is automatic.
\end{theorem}

\begin{example}[Adjacent transpositions]\label{ex:adj}
Let $c$ swap $2i{-}1\leftrightarrow 2i$ for every $i$ (so $c(n)=n+1$ if
$n$ is odd, $c(n)=n-1$ if $n$ is even). Here $g(n)=c(n)-n\in\{+1,-1\}$ is
bounded, yet $c$ is not eventually identity; Bob wins by local decoding.
This already shows the boundary is non-EI, not unbounded growth of~$g$.
\end{example}

\begin{example}[Block $3$-cycles on consecutive triples]\label{ex:cyc3}
Let $c$ cyclically rotate each triple $\{3i{-}2,3i{-}1,3i\}$, e.g.\
$3i{-}2\mapsto 3i{-}1\mapsto 3i\mapsto 3i{-}2$. This is locally decodable
(the bucketed pair of adjacent labels recovers $k_0$), and Bob wins.
\end{example}

Some non-EI permutations need a third probed label. The block $3$-cycle
written ``forwards'' is the canonical example.

\begin{definition}[Locally $K{=}3$ decodable]\label{def:k3}
$c$ is \emph{locally $3$-decodable} if the triple of first-\Yes{} times at
displacements $0,1,2$---equivalently the bucketed triple
$(c(k_0), c(k_0{+}1), c(k_0{+}2))$---determines~$k_0$.
\end{definition}

\begin{theorem}[$K{=}3$ winning theorem]\label{thm:k3}
If $c$ is locally $3$-decodable, Bob wins, via a period-$4$ triangle-wave
sweep of displacements $0,1,2,1,0,1,2,\dots$ with three detectors (first
\Yes{} at displacements $0$, $1$, $2$). The displacement stays in
$[0,2]$, so safety is automatic.
\end{theorem}

\begin{example}[Forward block $3$-cycle]\label{ex:blk}
Let $c(m)=m+1$ unless $3\mid m$, in which case $c(m)=m-2$ (so each block
$\{3b{+}1,3b{+}2,3b{+}3\}$ rotates forwards: $c(1)=2$, $c(2)=3$,
$c(3)=1$). Then $c(m)\le m+1$ always, so \emph{no} rightward probe
$[c(k_0{+}D) < T_1{+}D]$ ever separates the Phase-1 candidates---a
navigate-only strategy fails. But cell $3$ wakes early, at time~$2$ (since
$c(3)=1$), which splits $k_0=1$ from $k_0=2$ at displacement~$1$: at
$(D{=}1,\,t{=}3)$, the world $k_0=2$ stands on the already-awake cell~$3$
and reports \Yes, while $k_0=1$ stands on cell~$2$, still asleep
($c(2)=3\not<3$), and reports \No. The $K{=}3$ sweep, monitoring small
displacements from time~$0$, catches this early wake-up. So $c$ is
locally $3$-decodable, and Bob wins.
\end{example}

Example~\ref{ex:blk} is instructive: it is a concrete non-EI Bob-win
whose resolution \emph{requires} catching an early wake-up at a small
displacement, which is precisely the regime the local decoders are built
for---and precisely what naive ``oscillate then navigate right''
strategies miss.

% ------------------------------------------------------------
\subsection{Worked examples table}
\label{sec:table}

\begin{table}[h]
\centering
\renewcommand{\arraystretch}{1.25}
\begin{tabularx}{\textwidth}{@{}l l c X@{}}
\toprule
\textbf{Permutation $c$} & \textbf{EI?} & \textbf{Bob} & \textbf{Reason} \\
\midrule
Identity $c(n)=n$
  & yes & \textbf{loses}
  & $(2j,\,2j{+}1)$ collide: only even thresholds probeable
    (Rem.~\ref{rem:flat}). \\
Finitary (any finite-support perm.)
  & yes & \textbf{loses}
  & EI with $N=$ support bound; same collision for large $k_0$. \\
Adjacent transpositions
  & no & \textbf{wins}
  & locally $2$-decodable (Ex.~\ref{ex:adj}); bounded $g$. \\
Block $3$-cycle (rotated triples)
  & no & \textbf{wins}
  & locally $2$- or $3$-decodable (Ex.~\ref{ex:cyc3},\,\ref{ex:blk}). \\
Any $c$ rearranging arbitrarily large values
  & no & \textbf{wins}
  & not EI; distinguishable pairs (Thm.~\ref{thm:lemmaA}). \\
\bottomrule
\end{tabularx}
\caption{Outcomes by family. Eventually-identity $\iff$ Bob loses.}
\end{table}

% ============================================================
\section{Machine-verification status}
\label{sec:status}

The argument has been formalised in Lean~4 (with Mathlib) in the
development \texttt{griddles-p4-lean}. We report the status honestly: a
substantial core is machine-checked and axiom-clean, while a single
component---an explicit uniform winning strategy for the fully general
case---remains open. We do not overclaim.

For each item below, ``axiom-clean'' means the Lean \verb|#print axioms|
audit lists only the three standard foundational axioms
\verb|propext|, \verb|Classical.choice|, \verb|Quot.sound|, with
\emph{no} \verb|sorryAx| and no custom axioms. The entire library builds
(\verb|lake build| succeeds); the only \verb|sorry| in the whole
development is the one noted as open below.

\begin{table}[h]
\centering
\renewcommand{\arraystretch}{1.25}
\begin{tabularx}{\textwidth}{@{}l l >{\raggedright\arraybackslash}X@{}}
\toprule
\textbf{Component} & \textbf{Status} & \textbf{Lean location} \\
\midrule
Losing direction (Thm.~\ref{thm:losing})
  & \textbf{proven, axiom-clean}
  & \lean{Losing\dt EventuallyIdentity\us loses}; parity lock in
    \lean{Parity\dt lean}. \\
Distinguishability (Thm.~\ref{thm:lemmaA})
  & \textbf{proven, axiom-clean}
  & \lean{LemmaA\dt Indistinguishable\us implies\us eventually\us identity};
    crux \lean{indist\us Delta\us ge\us 2\us impossible}. \\
$K{=}2$ winning class (Thm.~\ref{thm:k2})
  & \textbf{proven, axiom-clean}
  & \lean{WinningRestricted\dt LocallyDecodable\us BobWins}; instances
    \lean{adjPerm\us BobWins}, \lean{cyc3Perm\us BobWins}. \\
$K{=}3$ winning class (Thm.~\ref{thm:k3})
  & \textbf{proven, axiom-clean}
  & \lean{WinningKProbe\dt LocallyK3Decodable\us BobWins}; instance
    \lean{blkPerm\us BobWins}. \\
Capstone (locally-decodable consistency)
  & \textbf{proven, axiom-clean}
  & \lean{Answer\dt locallyDecodable\us main}. \\
Main \emph{iff}, winning half, general $c$
  & \textbf{open}
  & \lean{Answer\dt main} ($\Leftarrow$ direction is a \lean{sorry}). \\
\bottomrule
\end{tabularx}
\caption{Formalisation status. Everything except the general winning
strategy is machine-checked and axiom-clean.}
\end{table}

\paragraph{What is proven.}
The losing direction is complete: eventually-identity~$c$ admits the
indistinguishable pair $(2j,2j{+}1)$, machine-verified through the
parallel-evolution lemma and the signal-match case analysis. The
distinguishability theorem is complete and axiom-clean, including its
crux for $\Delta\ge2$: the elementary forward-closure argument of
Section~\ref{sec:winning} (which earlier, by our own assessment, was
believed to need global surjectivity bookkeeping). The two
locally-decodable winning classes are complete, with the three concrete
instances $\texttt{adjPerm}$, $\texttt{cyc3Perm}$, $\texttt{blkPerm}$
each proven a Bob-win.

\paragraph{What is open---and why it is genuinely hard.}
The remaining gap is a single, explicit, uniformly-described strategy that
turns ``no pair is indistinguishable'' (Thm.~\ref{thm:lemmaA}) into a
formalised $\texttt{BobWins}$ for \emph{every} non-EI~$c$. The difficulty
is a real \emph{timing/resolution/reach} tension, not a missing lemma:

\begin{itemize}[topsep=2pt,itemsep=2pt]
  \item Bob occupies one displacement per tick, so he can finely monitor
  (period~$\approx 2$) only $O(1)$ displacements at once. Monitoring
  displacements $0,\dots,K$ simultaneously by a triangle wave costs time
  resolution $\approx 2K$, which can blur a discriminator whose
  separating window is $O(1)$.

  \item A discriminator that is an \emph{early} wake-up at displacement~$D$
  must be monitored finely \emph{from time~$0$} (it cannot be reached
  later: Example~\ref{ex:blk}).

  \item A discriminator that is a \emph{far} wake-up requires \emph{reaching}
  a large displacement, which precludes finely monitoring all small
  displacements meanwhile.
\end{itemize}

A single uniform strategy must reconcile these. We verified by direct
simulation that several natural candidate strategies are each incomplete
on a concrete family: a navigate-only ``monotone-right'' strategy fails on
the forward block $3$-cycle (Example~\ref{ex:blk}); ``$K{=}2$ or
navigate'' misses mid-range early discriminators of a block $5$-cycle; a
sequential staircase that probes displacement~$0$ first arrives at
displacement~$1$ too late to catch an early wake-up. The existence of
\emph{some} winning strategy for every non-EI~$c$ is confirmed by an
exhaustive belief-state (AND/OR) game solver with real safety, which
reports \textsc{lose} for eventually-identity inputs and \textsc{win}
from full ignorance for non-EI inputs, stably as the search horizon
grows. But a single explicit, provably-complete, formalisable
construction is an intricate open problem---plausibly the intended hard
core of the puzzle.

\paragraph{Summary.}
The \emph{answer} (Theorem~\ref{thm:main}) is settled, with the losing
direction and the distinguishability theorem machine-checked and
axiom-clean, and explicit machine-checked winning strategies for natural
families. The one component not yet formalised is a uniform winning
strategy for arbitrary non-EI~$c$; its existence is established by an
exhaustive game solver, while an explicit formalisable strategy remains
open.

% ============================================================
\section*{Acknowledgements}
\addcontentsline{toc}{section}{Acknowledgements}
% ============================================================

The solution to this puzzle --- including the characterisation of Bob's
winning condition, the belief-state strategy construction, and the
fixed-$K$ impossibility argument --- was developed in collaboration with
\href{https://claude.ai}{Claude} (Anthropic).  The simulation and
belief-state game-solver code, the Lean~4 formalisation, and this writeup
were subsequently produced in collaboration with
\href{https://claude.ai/code}{Claude Code} (Anthropic).

% ============================================================

\end{document}
